\chapter{Expanding the variance}\label{ch:variance}

We open the square in the definition of $V_f$ and separate three kinds of terms: a diagonal term, a mixed term involving the densities, and an off-diagonal term that counts pairs of squarefree numbers whose difference is a multiple of the modulus. The moduli $f(k)$ with $k$ below a cutoff $z$ are removed by a known upper bound. (Archived source: proposition P2.)

\section{Notation}

For real $x\ge2$ and integers $1\le m\le x$ put
\begin{equation}\label{eq:Cx}
C_x(m)=\sum_{1\le n\le x-m}\mu^2(n)\mu^2(n+m).
\end{equation}
For a modulus $q$ write $H(q)=\sum_{l=1}^qg(q,l)^2$. For real $1\le z\le y$ with $Z=f(z)\le Y=f(y)\le x$ define
\begin{align}
V_f(x;z,y)&=\sum_{z<k\le y}f'(k)\sum_{l=1}^{f(k)}E(x;f(k),l)^2,\label{eq:Vzy}\\
T_f(z,y)&=\sum_{z<k\le y}f'(k)\sum_{l=1}^{f(k)}g(f(k),l)^2=\sum_{z<k\le y}f'(k)H(f(k)),\label{eq:Tf}\\
S_f(x;z,y)&=\sum_{z<k\le y}f'(k)\sum_{1\le h\le x/f(k)}C_x(hf(k)).\label{eq:Sf}
\end{align}
Since $f$ is increasing, $f(k)\le Y\le x$ for $k\le y$, and every argument $hf(k)$ of $C_x$ in \eqref{eq:Sf} lies in $[1,x]$. Note that $V_f(x,y)=V_f(x,z)+V_f(x;z,y)$ for $1\le z\le y$.

\section{Counting pairs in a residue class}

\begin{lemma}\label{lem:pairs}
For real $x\ge2$ and every integer $1\le q\le x$,
\[
\sum_{l=1}^qQ(x;q,l)^2=\sum_{n\le x}\mu^2(n)+2\sum_{1\le h\le x/q}C_x(hq).
\]
\end{lemma}

\begin{proof}
The left side counts ordered pairs $(n_1,n_2)$ of squarefree integers in $[1,x]$ with $n_1\equiv n_2\bmod q$. Pairs with $n_1=n_2$ give the first term. For the others, the positive difference is $hq$ for exactly one integer $1\le h\le x/q$; the pairs with $n_2-n_1=hq$ are counted by $C_x(hq)$, and reversing the order gives the factor two.
\end{proof}

\begin{lemma}\label{lem:mixed}
For every integer $q\ge1$ and real $x\ge1$, the mixed sum
\[
\Delta(x;q)=\sum_{l=1}^qg(q,l)E(x;q,l)
\]
satisfies $|\Delta(x;q)|\le3\sqrt x\,\tau(q)/q$.
\end{lemma}

\begin{proof}
Put $q_1=\prod_{p\parallel q}p$ and $A_q=\prod_{p\mid q}(1-p^{-2})^{-1}$. By \eqref{eq:c-euler}, $cA_q=\prod_{p\nmid q}(1-p^{-2})\le1$. For squarefree $n$ the zero case of \cref{prop:density} cannot occur, and expanding $\prod_{p\mid(q_1,n)}(1-p^{-1})$ in \eqref{eq:density-closed} gives
\[
\mu^2(n)g(q,n)=\frac{cA_q}q\mu^2(n)\sum_{a\mid q_1,\ a\mid n}\frac{\mu(a)}a.
\]
The squarefree $n\le x$ divisible by $a$ are counted by $Q(x;a,0)$, and \cref{prop:density} gives $|Q(x;a,0)-x\,g(a,0)|\le3\sqrt x$. Hence, with $B(q)=\frac{cA_q}q\sum_{a\mid q_1}\frac{\mu(a)}a\,g(a,0)$,
\begin{equation}\label{eq:mixed-B}
\Bigl|\sum_{n\le x}\mu^2(n)g(q,n)-xB(q)\Bigr|\le\frac{3\sqrt x}q\sum_{a\mid q_1}\frac1a\le\frac{3\sqrt x\,\tau(q)}q.
\end{equation}
Since $g(q,n)$ depends only on $n\bmod q$,
$\sum_{n\le t}\mu^2(n)g(q,n)=\sum_{l=1}^qg(q,l)Q(t;q,l)$. Dividing by $t$ and letting $t\to\infty$, the right side tends to $H(q)$ and, by \eqref{eq:mixed-B}, the left side tends to $B(q)$. Thus $B(q)=H(q)$, and since $\Delta(x;q)=\sum_{l}g(q,l)Q(x;q,l)-xH(q)=\sum_{n\le x}\mu^2(n)g(q,n)-xH(q)$, the claim follows from \eqref{eq:mixed-B}.
\end{proof}

\section{The variance on an interval of moduli}

\begin{proposition}\label{prop:reduction}
For every $\eps>0$ there is $C(f,\eps)>0$ such that, uniformly for real $x\ge2$ and $1\le z\le y$ with $Z=f(z)\le Y=f(y)\le x$,
\[
\Bigl|V_f(x;z,y)-\bigl[cx(Y-Z)+2S_f(x;z,y)-x^2T_f(z,y)\bigr]\Bigr|\le C(f,\eps)\bigl(x^{3/2+\eps}+xY^{1-1/d}\bigr).
\]
\end{proposition}

\begin{proof}
Fix $q=f(k)\le x$. Since $Q=E+xg$, expanding $\sum_lQ^2$ and using \cref{lem:pairs} and the definition of $\Delta$ gives the exact identity
\[
\sum_{l=1}^qE(x;q,l)^2=\sum_{n\le x}\mu^2(n)+2\sum_{1\le h\le x/q}C_x(hq)-x^2H(q)-2x\Delta(x;q).
\]
Multiply by $f'(k)$ and sum over $z<k\le y$. With $\mathcal W=\sum_{z<k\le y}f'(k)$, the difference to be estimated is exactly
\[
\mathcal R=\Bigl(\sum_{n\le x}\mu^2(n)-cx\Bigr)\mathcal W+cx\bigl(\mathcal W-(Y-Z)\bigr)-2x\sum_{z<k\le y}f'(k)\Delta(x;f(k)).
\]
\emph{Diagonal error.} \Cref{prop:density} with $q=1$ gives $|\sum_{n\le x}\mu^2(n)-cx|\le3\sqrt x$, since $g(1,0)=\sum_m\mu(m)m^{-2}=c$. By \eqref{eq:weight-count} and $Y\le x$, the first term is at most $3(1+K_1)\sqrt xY\le3(1+K_1)x^{3/2}$.

\emph{Endpoint error.} By \eqref{eq:weight-count}, the second term is at most $cK_1xY^{1-1/d}$.

\emph{Mixed error.} Let $\eta=\eps/2$. By \cref{lem:mixed}, \cref{lem:divisor} and $f(k)\le x$, each term satisfies $|f'(k)\Delta(x;f(k))|\le3D_\eta x^{1/2+\eta}f'(k)/f(k)$. Summing with \eqref{eq:sum-w} and using $\log2x\le J_\eta x^\eta$ (\cref{lem:sums}(d)), we find that the third term is at most
\[
2x\cdot3D_\eta x^{1/2+\eta}\cdot2U\log(2x)\le12D_\eta UJ_\eta\,x^{3/2+\eps}.
\]
Adding the three bounds proves the proposition with $C(f,\eps)=3(1+K_1)+12D_{\eps/2}UJ_{\eps/2}+cK_1$. The argument also covers $z=y$ and intervals containing no integer.
\end{proof}

\section{The initial moduli}

For the moduli $f(k)$ with $k\le z$ we use the full-variance bound of Brüdern, Granville, Perelli, Vaughan and Wooley. In their notation $\mu_k$ is the indicator of the $k$-free numbers, so $\mu_2=\mu^2$, and their parameter $k$ is unrelated to our summation variable.

\begin{quoted}[{\cite[Lemma~2.2 and (2.6)--(2.7)]{bgpvw1998}}]\label{quoted:bgpvw}
For positive integers $a,q$ define $E_k(x;q,a)$ by
\[
\sum_{\substack{1\le n\le x\\n\equiv a\ (\mathrm{mod}\ q)}}\mu_k(n)=\frac x{\zeta(k)}g_k(q,(q,a))+E_k(x;q,a),
\qquad
\Upsilon_k(x,Q)=\sum_{1\le q\le Q}\sum_{a=1}^q|E_k(x;q,a)|^2.
\]
Suppose that $Q$ and $x$ are positive real numbers and that $k\ge2$ is an integer. Then for each $\eps>0$,
\[
\Upsilon_k(x,Q)\ll x^{2/k+\eps}Q^{2-2/k}\quad(1\le Q\le x),\qquad\Upsilon_k(x,Q)\ll Q^2\log(2Q)\quad(Q>x).
\]
\end{quoted}

The inner sum in $\Upsilon_k$ runs over all residue classes, coprime to $q$ or not. We check that the normalisation agrees with ours without using the formula for $g_k$. Fix $1\le a\le q$ and apply the lemma with $k=2$, $Q=q$, $\eps=\frac12$ and $x=t\ge\max\{2,q\}$. By positivity, $|E_2(t;q,a)|^2\le\Upsilon_2(t,q)\ll t^{3/2}q$, so $E_2(t;q,a)/t\to0$ as $t\to\infty$. Dividing the definition by $t$ therefore shows that $\zeta(2)^{-1}g_2(q,(q,a))=\lim_tQ(t;q,a)/t=g(q,a)$. Hence $E_2=E$, and with $k=2$ the lemma states that
\begin{equation}\label{eq:full-variance}
V(x,Q)=\sum_{q\le Q}\sum_{l=1}^qE(x;q,l)^2\le B_\eps x^{1+\eps}Q\qquad(1\le Q\le x)
\end{equation}
for every $\eps>0$, where $B_\eps$ depends only on $\eps$. The same bound is recalled as (1.5) in \cite{vaughan2005}.

\begin{proposition}\label{prop:initial}
For every $\eps>0$ there is $C(f,\eps)$ such that, uniformly for real $x\ge2$ and real $z\ge1$ with $f(z)\le x$,
\[
0\le V_f(x,z)\le C(f,\eps)\,x^{1+\eps}f(z)^{2-1/d}.
\]
\end{proposition}

\begin{proof}
Let $Z=f(z)$. By \cref{lem:poly}(a) the moduli $f(k)$, $k\le z$, are distinct integers in $[2,Z]$, and by \cref{lem:poly}(e) their weights satisfy $0<f'(k)\le\kappa_2\kappa_1^{-\alpha}Z^\alpha$ with $\alpha=1-1/d$. Each inner sum $\sum_lE^2$ is nonnegative. We may therefore first bound each weight by this maximum and then enlarge the set of moduli to all $q\le Z$:
\[
0\le V_f(x,z)\le\kappa_2\kappa_1^{-\alpha}Z^{\alpha}\sum_{q\le Z}\sum_{l=1}^qE(x;q,l)^2\le\kappa_2\kappa_1^{-\alpha}B_\eps\,x^{1+\eps}Z^{1+\alpha}.
\]
Here \eqref{eq:full-variance} is applicable because $1\le2\le Z\le x$, and $1+\alpha=2-1/d$.
\end{proof}

\begin{remark}
It is essential to use positivity before replacing the weights by their maximum. Bounding a sparse sum by the full sum over all $q\le Z$ costs the factor $Z^{1-1/d}$. This loss is one of the two sources of the exponent $1/(4d-1)$ in \cref{thm:main}.
\end{remark}
